% RESULTADO_RECUPERADO desde II REV09. Véase MAPA_PROCEDENCIA.json.
% Selección de líneas y cambios mecánicos explícitos; ninguna prueba nueva.
\subsubsection{Carácter de acción y procedencia de sus coeficientes}
\label{v:ant:sec:revision-planck}

La constante de Planck reducida, \(\hbar\), es el objeto físico que
motiva la realización de la sección de acción angular. Su posición en
la dependencia electrónica debe especificarse en dos niveles. La norma
determinantal \(\Sigma_H\) determina la década de la escala. El carácter
\(H_5\) y el retorno regional determinan después las coordenadas de las
secciones de acción. Esta separación permite exponer el resultado
necesario para el electrón sin anticipar los desarrollos completos de
otros sectores de la mecánica dimensional.

La procedencia del carácter de quinto orden puede explicitarse mediante
invariantes ya fijados en el soporte. En la coordenatización central se considera
\[
 B_c=\begin{pmatrix}7&2\\2&7\end{pmatrix},
 \qquad \lambda_+=9,\quad\lambda_-=5.
\]
El calendario tiene longitud \(12\), la órbita del paso \(3\) tiene
longitud \(4\), y el cociente censal pertinente es
\(104976/1944=54\). Los coeficientes del carácter satisfacen
\begin{equation}
 \frac9{16}=\frac{\lambda_+}{4^2},\quad
 \frac59=\frac{\lambda_-}{\lambda_+},\quad
 \frac7{48}=\frac{\operatorname{tr}(B_c)/2}{4\cdot12},\quad
 \frac1{54}=\frac{1944}{104976}.
 \label{v:ant:eq:revision-h5-coeficientes}
\end{equation}
La asignación de estos invariantes a los grados \(2,3,4,5\), con la
alternancia de signos de \eqref{v:ant:eq:el-h5}, pertenece a la definición del
carácter analítico utilizado. Explicitarla es esencial: el conjunto de
cardinales, considerado sin el carácter y sin el orden de sus grados,
admitiría otras expresiones. El contenido constructivo reside en la
composición del soporte, la orientación y la regla analítica declarada.

\subsubsection{Continuación regional y ecuación de renovación}

El selector \(169\) determina el impulso de grado seis. La prolongación
posterior se organiza por la transferencia \(D_A\) y una normalización
unitaria de la línea determinantal. Estas reglas admiten una formulación
en series formales antes de sustituir \(z=\alpha\), con \(D_A\) ya
construido. Escribamos
\[
 \mathscr C(z)=169z^6+R(z),\qquad R(z)\in z^7\mathbb R[[z]].
\]
La concatenación de una prolongación aumenta el grado en una unidad y
multiplica su peso por \(D_A\). La primera prolongación estricta tiene
peso \(D_A\). Por tanto, su ecuación de renovación es
\begin{equation}
 R(z)=D_Az^7+D_AzR(z).
 \label{v:ant:eq:revision-renovacion}
\end{equation}

\begin{proposition}[Unicidad formal de la continuación normalizada]
La ecuación \eqref{v:ant:eq:revision-renovacion} determina una única serie y
\[
 \mathscr C(z)=169z^6+\frac{D_Az^7}{1-D_Az}
             =169z^6+\sum_{n=7}^{\infty}D_A^{n-6}z^n.
\]
La realización es analítica en \(|D_Az|<1\).
\end{proposition}
\begin{proof}
El elemento \(1-D_Az\) tiene término independiente uno y es invertible
en el anillo de series formales. Resolver la ecuación produce la serie
exhibida. Equivalentemente, el coeficiente de grado siete es \(D_A\)
y cada coeficiente posterior es \(D_A\) veces el anterior. La serie
geométrica converge absolutamente en el disco indicado.
\end{proof}

La normalización de la primera prolongación es un dato operativo
indispensable. Si se conservaran únicamente el impulso de grado seis y
la razón de concatenación, las series
\[
 169z^6+\lambda\frac{D_Az^7}{1-D_Az}
\]
seguirían compartiendo ambos datos para cualquier \(\lambda\).
La regla unitaria fija \(\lambda=1\) antes de efectuar la evaluación.
Así queda localizada la contribución de cada regla a la unicidad del
lector, y la expresión racional conserva la continuación completa.

\subsubsection{Constante de Planck reducida y retorno de la sección de acción}

En la coordenatización dimensional de \eqref{v:ant:eq:el-secciones-accion}, la sección
anterior al retorno es la construcción del modelo destinada al contraste
con la constante de Planck reducida:
\[
 \hbar_{\mathrm{pre}}=H_5\,10^{-34}\mathcal U_S.
\]
La sección posterior conserva, además, la compensación regional
\(90\mathscr C(\alpha)/\pi\). Ambas pertenecen a la misma recta de
acción y su cociente \(R_{\mathrm{act}}\) mide el efecto multiplicativo
del retorno. El paso a \(h\) mediante \(2\pi\) utiliza la relación
entre la parametrización angular y una vuelta completa.
